\documentclass[12pt, twoside]{article}
\input{../preamble}

\fancyhead[LO]{La Scuola d'Italia / Dr.\ Huson / IB Math: Practice \\* 15 April 2026}
\fancyhead[RO]{\fbox{\begin{minipage}{6cm}\footnotesize First \& Last name:\\[0.5cm]\end{minipage}}}
\fancyfoot[C]{\thepage}

\begin{document}
\subsubsection*{11.12 Classwork: Exponential Functions and Compound Interest (calculator)}

\begin{enumerate}[itemsep=0.15cm]

%--- Problem 1: Car depreciation + compound interest [6 marks] ---
%--- Source: AA SL 2023 May P2 TZ1 Q02 ---
\item[1.] [Maximum mark: 6]

The value of a car is given by the function $C = 40\,000(0.91)^t$,
where $t$ is in years since 1 January 2023 and $C$ is in USD (\$).

\begin{enumerate}[itemsep=0.15cm]
  \item Write down the annual rate of depreciation of the car. \hfill [1]
  \item Find the value of the car on 1 January 2028. \hfill [2]
\end{enumerate}

Alvie wants to buy this car. On 1 January 2023, he invested \$15\,000
in an account that earns 3\,\% annual interest compounded yearly.
He makes no further deposits to, or withdrawals from, the account.

Alvie wishes to buy this car for its value on 1 January 2028.
In addition to the money in his account, he will need an extra \$$M$.

\begin{enumerate}[itemsep=0.15cm]
  \item[(c)] Find the value of $M$. \hfill [3]
\end{enumerate}

\begin{tikzpicture}
  \draw (-0.6,0) rectangle (11,13);
  \foreach \y in {1,2,...,12} \draw[dotted] (0,\y)--(10.5,\y);
\end{tikzpicture}

\vspace{0.2cm}

\newpage

%--- Problem 2: Compound interest + geometric series [16 marks] ---
%--- Source: AA SL 2021 May P2 TZ1 Q07 ---
Answer on a separate sheet of lined paper. Label each section clearly and leave a blank line between each part.
\item[2.] [Maximum mark: 16]

Two friends Amelia and Bill, each set themselves a target of saving \$20\,000.
They each have \$9000 to invest.

\begin{enumerate}[itemsep=0.15cm]
  \item Amelia invests her \$9000 in an account that offers an interest rate
  of 7\,\% per annum compounded annually.

  \begin{enumerate}[itemsep=0.1cm]
    \item[(i)] Find the value of Amelia's investment after 5 years to
    the nearest hundred dollars. \hfill [2]
    \item[(ii)] Determine the number of years required for Amelia's
    investment to reach the target. \hfill [3]
  \end{enumerate}

  \item Bill invests his \$9000 in an account that offers an interest rate
  of $r$\,\% per annum compounded monthly, where $r$ is set to two
  decimal places.

  Find the minimum value of $r$ needed for Bill to reach the target
  after 10 years. \hfill [3]

  \item A third friend Chris also wants to reach the \$20\,000 target.
  He puts his money in a safe where he does not earn any interest.
  His system is to add more money to this safe each year.
  Each year he will add half the amount added in the previous year.

  \begin{enumerate}[itemsep=0.1cm]
    \item[(i)] Show that Chris will never reach the target if his initial
    deposit is \$9000. \hfill [4]
    \item[(ii)] Find the amount Chris needs to deposit initially in order
    to reach the target after 5 years. Give your answer to the nearest
    dollar. \hfill [4]
  \end{enumerate}
\end{enumerate}

\end{enumerate}
\end{document}
